\section{Heavy-Light Decomposition}
\whenuse{Updates o consultas sobre caminos y subárboles, combinando el árbol
con Fenwick o Segment Tree.}

Para cada vértice se elige como hijo pesado al hijo con mayor subárbol. Las
aristas pesadas forman cadenas cuyos vértices ocupan posiciones consecutivas.
Así, un camino se divide en $O(\log n)$ intervalos del arreglo base.

\complexity{Build $O(n)$; camino con Segment Tree $O(\log^2 n)$; subárbol
$O(\log n)$; memoria $O(n)$.}

\begin{lstlisting}[style=cpp]
struct HLD {
    int n,timer=0;
    vector<vector<int>> tree;
    vector<int> parent,depth,size,heavy,head,pos;
    HLD(vector<vector<int>> tree,int root=0):n(tree.size()),
    tree(move(tree)),parent(n,-1),depth(n),size(n),heavy(n,-1),
        head(n),pos(n) {
        assert(n>0);
        dfs_size(root);
        dfs_decompose(root,root);
    }
    int dfs_size(int v) {
        size[v]=1;
        int largest=0;
        for(int to:tree[v]) if(to!=parent[v]) {
            parent[to]=v;
            depth[to]=depth[v]+1;
            int child_size=dfs_size(to);
            size[v]+=child_size;
            if(child_size>largest) largest=child_size,heavy[v]=to;
        }
        return size[v];
    }
    void dfs_decompose(int v,int chain_head) {
        head[v]=chain_head;
        pos[v]=timer++;
        if(heavy[v]!=-1) dfs_decompose(heavy[v],chain_head);
        for(int to:tree[v])
            if(to!=parent[v]&&to!=heavy[v]) dfs_decompose(to,to);
    }
    // Intervalos [l,r) que cubren el camino a--b.
    template<class Callback>
    void for_each_path(int a,int b,Callback callback) const {
        while(head[a]!=head[b]) {
            if(depth[head[a]]<depth[head[b]]) swap(a,b);
            callback(pos[head[a]],pos[a]+1);
            a=parent[head[a]];
        }
        if(depth[a]>depth[b]) swap(a,b);
        callback(pos[a],pos[b]+1);
    }
    // Los vertices del subarbol de v son consecutivos.
    pair<int,int> subtree_interval(int v) const {
        return {pos[v],pos[v]+size[v]};
    }
};
\end{lstlisting}

Ejemplo con un Segment Tree de suma:

\begin{lstlisting}[style=cpp]
HLD hld(tree,0);
vector<ll> base(n);
for(int v=0;v<n;++v) base[hld.pos[v]]=value[v];
SegmentTree segment(base);

// Cambiar el valor del vertice v.
segment.set_value(hld.pos[v],new_value);

// Suma de los vertices del camino a--b.
ll path_sum=0;
hld.for_each_path(a,b,[&](int l,int r) {
    path_sum+=segment.query(l,r);
});

// Suma del subarbol de v.
auto [l,r]=hld.subtree_interval(v);
ll subtree_sum=segment.query(l,r);
\end{lstlisting}

Para guardar valores en aristas, almacenar cada arista en la posición de su
vértice hijo. Para excluir el LCA de la consulta:

\begin{lstlisting}[style=cpp]
hld.for_each_path(a,b,[&](int l,int r) {
    // Adaptar para excluir la posicion del LCA.
});
\end{lstlisting}

\warning{El callback no entrega necesariamente los segmentos en el orden del
camino. Esto no importa para suma, mínimo, máximo o XOR; para operaciones no
conmutativas se deben separar y orientar los segmentos de ambos lados.}